\subsection{Simulator and Training Details} \label{sup:sim}

The design of the reward structure and termination conditions plays a critical role in shaping the learned policy. Overly strict rewards can prevent learning altogether, while rewards that are too lenient often lead to unsafe or suboptimal behavior. In practice, tuning the reward function is crucial for the success of reinforcement learning algorithms \cite{Silver2021}.

We employ a shaped reward function composed of five components:

\begin{itemize}
    \item A constant survival reward \( C_r \) to prevent the agent from learning to terminate episodes early.
    \item Penalties for deviations from the desired behavior:
        \begin{itemize}
            \item Position error penalty \( C_p \)
            \item Velocity error penalty \( C_v \)
            \item Orientation error penalty \( C_q \)
            \item Action magnitude penalty \( C_a \)
        \end{itemize}
\end{itemize}

To encourage stable and progressively more precise control, a reward curriculum is applied by linearly increasing the penalties during training. The initial and final values are summarized in Table~\ref{tab:learning_types}.

\begin{table}[h]
    \centering
    \begin{tabular}{lccccc}
        \toprule
        & \( C_p \) & \( C_v \) & \( C_a \) & \( C_q \) & \( C_r \) \\
        \midrule
        Start & 1.0 & 0.005 & 0.01 & 0.25 & 1.0 \\
        End   & 2.0 & 0.2   & 0.03 & 0.25 & 1.0 \\
        \bottomrule
    \end{tabular}
    \caption{Reward parameters used during training. Penalties increase linearly across epochs, updated every 10 epochs.}
    \label{tab:learning_types}
\end{table}

\paragraph{Drone Dynamics.}
The simulated drone uses a simple physics model. Motor thrust is derived from RPM via a second-order polynomial:
\[
T = c_0 + c_1 \cdot \text{rpm} + c_2 \cdot \text{rpm}^2
\]
The motor dynamics are modeled using a first-order low-pass filter:
\[
\Delta \text{rpm} = \frac{\text{rpm}_{\text{des}} - \text{rpm}_{\text{curr}}}{\tau}
\]
where \( \tau \) represents the motor time constant. Body-frame dynamics are integrated and then transformed to the world frame.

\paragraph{Training Hyperparameters.}
Table~\ref{tab:hyperparameters} provides an overview of the hyperparameters used across all methods. Common parameters are listed first, followed by method-specific values.

\begin{table}[htbp]
\centering
\caption{Training hyperparameters for all methods.}
\label{tab:hyperparameters}
\begin{tabular}{lllll}
\toprule
\textbf{Parameter} & \textbf{BC} & \textbf{TD3BC} & \textbf{TD3BC Online} & \textbf{TD3} \\
\midrule
\multicolumn{5}{l}{\textbf{Common Parameters}} \\
\midrule
Hidden sizes        & [256, 128] & [256, 256] & [256, 128] & [256, 128] \\
Learning rate       & 1e-3       & 1e-3       & 1e-3       & 1e-3       \\
Batch size          & 50         & 125        & 350        & 256        \\
Buffer size         & 1M         & 1M         & 1M         & 1M         \\
Slope               & 2          & 2          & 2          & 2          \\
Slope schedule      & adaptive   & adaptive   & adaptive   & adaptive   \\
Scheduling order    & 3          & 3          & 3          & 3          \\
\midrule
\multicolumn{5}{l}{\textbf{Method-Specific Parameters}} \\
\midrule
Warm-up steps       & 50         & 50         & 50         & --         \\
Policy noise        & 0.0        & 0.0        & 0.2        & --         \\
Noise clip          & --         & 0.5        & 0.5        & --         \\
\(\alpha\) (TD3BC)  & --         & 2.5        & 2.5        & 0.0        \\
\(\tau\) (Target)   & --         & 0.005      & 0.005      & 0.005      \\
\(\gamma\) (Discount) & --       & 0.99       & 0.99       & 0.99       \\
BC coefficient      & --         & --         & 0.2        & --         \\
BC decay factor     & --         & --         & 0.99       & --         \\
Training epochs     & 300        & 300        & 300        & 300        \\
Steps per epoch     & --         & --         & 10000      & 5000       \\
Test episodes       & 20         & 20         & 20         & 20         \\
\bottomrule
\end{tabular}
\end{table}
